\documentclass[10pt,a4paper]{amsart}
\usepackage[margin=19mm]{geometry}
\usepackage{amsmath,amssymb,mathtools}
\usepackage[scr=rsfs,cal=euler]{mathalfa}
\usepackage{xfrac}
\usepackage[unicode,hidelinks,bookmarks=false]{hyperref}
\newcommand{\ie}{i.e.\ }
\newcommand{\cf}{cf.\ }
\newcommand{\resp}{respectively\ }
\newcommand{\Subsection}{\subsection}
\providecommand{\nobreakdash}{\nobreak\hbox{-}}
\setlength{\emergencystretch}{2em}
\multlinegap0pt
\usepackage{bm}
\usepackage{bbm}
\usepackage{dsfont}
\usepackage{xspace}
\usepackage{cancel}
\usepackage{mathtools}
\usepackage{amsthm}
\makeatletter
\@ifundefined{theorem}{\newtheorem{theorem}{Theorem}}{}
\@ifundefined{lemma}{\newtheorem{lemma}[theorem]{Lemma}}{}
\makeatother
\title[Anisotropic mean curvature flow with contact angle and Neuma]{Anisotropic mean curvature flow with contact angle and Neumann boundary conditions in arbitrary dimensions}
\author{Can Cui, Nung Kwan Yip}
\date{Source: arXiv:2510.22146v1 (CC BY 4.0)}
\begin{document}
\maketitle

\noindent\emph{Source and licence.} Anisotropic mean curvature flow with contact angle and Neumann boundary conditions in arbitrary dimensions. Version arXiv:2510.22146v1 (CC BY 4.0), distributed under the Creative Commons Attribution 4.0 International licence, as recorded in the arXiv licence metadata for this exact version and shown on the abstract page https://arxiv.org/abs/2510.22146v1. Full rights notice: licenses/arxiv-2510.22146-CC-BY-4.0.txt.\par\smallskip

\noindent\textit{Editorial scope.} Counted unit: the Lemma whose source label is \texttt{lem:2.4} in the article (source0.tex lines 361--379, source label \texttt{lem:2.4}), delivered together with the article's own complete original proof of it (source0.tex lines 380--456, one \texttt{proof} environment of 77 lines). No mathematics is written, repaired or re-derived: statement and proof are the article's own text line for line. Required context retained uncounted: homc3 (lines 211--213); ONB (lines 217--219); lem:2.1 (lines 263--268); lem:2.2 (lines 299--304); lem:2.3 (lines 326--331). Every definition, display and label of those lines is retained unchanged. Omitted from this package and not redistributed: 1--210, 214--216, 220--262, 269--298, 305--325, 332--360, 457--1546. Nothing inside a retained excerpt is omitted or altered: every retained body is delivered exactly as the article's lines are written, so no omission receipt of any kind accompanies this package. Citations: the retained excerpts carry no citation command at all, so no bibliography entry of the article is transcribed and no reference is redistributed. Reference adaptation: none was needed and none was made; every reference inside the delivered unit resolves inside it, so no unresolved reference or citation is produced. Printed slips kept literal: no printed word, symbol, hypothesis or number of the counted statement or of its proof was altered. Layout and class: the article's own class is \texttt{article} with packages including graphicx, appendix, amsmath, amssymb, amsthm, amsrefs, enumitem, xcolor, tikz, hyperref, lineno, geometry, endnotes; the wrapper is the lane's pinned \texttt{amsart} editorial wrapper at 10pt on A4, which restates the theorem-like environments the retained text needs and adds the article's own macro definitions that the retained text needs (none), with extra wrapper packages (bm, bbm, dsfont, xspace, cancel, mathtools). The delivered numbering is the unit's own. Assets: no retained span contains a figure environment, a graphics-inclusion command, a PDF-page inclusion, a table or a TikZ picture; no image file is copied into this package, the package declares no asset dependency and no image file is redistributed with it. Licence: version arXiv:2510.22146v1 of this article is distributed under the Creative Commons Attribution 4.0 International licence, as recorded in the arXiv licence metadata for this exact version and shown on the abstract page https://arxiv.org/abs/2510.22146v1. A full rights notice is delivered at licenses/arxiv-2510.22146-CC-BY-4.0.txt. Characters: every retained excerpt is pure ASCII, so the delivered engine, XeLaTeX with its default Latin Modern text font, reports no missing character.
\begin{equation} \label{homc3}
    D^3 F\Big|_p (p,\cdot,\cdot) = -D^2 F\Big|_p (\cdot,\cdot).
\end{equation}

\begin{equation}\label{ONB}
\{\phi^1, \cdots, \phi^n\}\,\,\,\text{with $\phi^n = \frac{p'}{|p'|}$ for $p'\neq0$.}
\end{equation}

\begin{lemma} \label{lem:2.1}
    There exists a positive constant $C_1$ such that for any $q'\in \mathbb{R}^n$ with $|q'|=1$, we have
\begin{equation}\label{eq:lem2.1}
    \left|D^2 F\Big|_{(p',-1)} ((p',0), (q',0))\right| \leq \frac{C_1}{|p'|^2}.
\end{equation}
\end{lemma}

\begin{lemma} \label{lem:2.2}
There exists a positive constant $C_2$ such that for any $q',r'\in \mathbb{R}^n$ with $|q'|=|r'|=1$, we have
\begin{equation}
  \left|  D^3 F\Big|_{(p',-1)} ((0,1),(q',0),(r',0) ) \right| \leq \frac{C_2}{|p'|^3}.
\end{equation}
\end{lemma}

\begin{lemma} \label{lem:2.3}
There exists a positive constant $C_3$ such that 
\begin{equation}
  \left|  D^3 F\Big|_{(p',-1)} ((0,1),(0,1),(0,1) \right| \leq \frac{C_3}{|p'|^3}.
\end{equation}
\end{lemma}

\begin{lemma} \label{lem:2.4}
Recall the basis from \eqref{ONB} with $\phi^n = \frac{p'}{|p'|}$. There exist positive constants $C_4, C_5, C_6$ depending only on $F$ such that for any $1\leq \alpha,\beta \leq n-1$, we have 
\begin{equation*}
    |T^3 (\phi^{\alpha}, \phi^{\alpha}, \phi^{\beta})| \leq \frac{C_4}{|p|},
\end{equation*}
\begin{equation*}
     |T^3 (\phi^{\alpha}, \phi^{\alpha}, \phi^{n})|,\ |T^3 (\phi^{\alpha}, \phi^n, \phi^n)|, \  |T^3 (\phi^{n}, \phi^{n}, \phi^{\alpha})| \leq \frac{C_4}{|p|^3},
\end{equation*}
\begin{equation*}
    | T^3 (\phi^{\alpha}, \phi^n, \phi^{\beta})| \leq \frac{C_4}{|p|} \quad (\alpha\neq \beta),
\end{equation*}
\begin{equation*}
     \frac{C_5}{|p|} \leq T^3 (\phi^{\alpha}, \phi^n, \phi^{\alpha}) \leq \frac{C_6}{|p|},
\end{equation*}
\begin{equation*}
\frac{C_5}{|p|^3} \leq T^3 (\phi^{n}, \phi^{n}, \phi^n)  \leq \frac{C_6}{|p|^3}.
\end{equation*}
%and $\alpha\neq\beta$.
\end{lemma}

\begin{proof}
Firstly, from \eqref{homc3}, we calculate the following to apply the results from Lemma \ref{lem:2.1} -- \ref{lem:2.3} to $D^3 F$ evaluated on basis elements,
\begin{equation*}
\begin{split}
     & D^3 F\Big|_{(p',-1)} ((p',0),(q',0),(r',0) ) \\
     = & - D^2 F\Big|_{(p',-1)} ((q',0), (r',0)) + D^3 F\Big|_{(p',-1)} ((0,1),(q',0),(r',0) ). \\
\end{split}
\end{equation*}
In addition, we have
\begin{equation*}
\begin{split}    
   &  D^3 F\Big|_{(p',-1)} ((p', 0), (p', 0), (p', 0)) \\
    = & - D^2 F\Big|_{(p',-1)} ((p', 0), (p', 0)) + D^3 F\Big|_{(p',-1)} ((0, 1), (p', 0), (p', 0)) \\
    & D^3 F\Big|_{(p',-1)} ((0, 1), (p', 0), (p', 0))  \\
    = & D^3 F\Big|_{(p',-1)} ((0, 1), (p', -1), (p', -1)) + D^3 F\Big|_{(p',-1)} ((0, 1), (p', -1), (0, 1)) \\
    & + D^3 F\Big|_{(p',-1)} ((0, 1), (0, 1), (p', -1)) + D^3 F\Big|_{(p',-1)} ((0, 1), (0, 1), (0, 1)) \\
    = & -2 D^2 F\Big|_{(p',-1)} ((0, 1), (0, 1)) + D^3 F\Big|_{(p',-1)} ((0, 1), (0, 1), (0, 1)) \\
    = & -2 D^2 F\Big|_{(p',-1)} ((p', 0), (p', 0)) + D^3 F\Big|_{(p',-1)} ((0, 1), (0, 1), (0, 1)).
\end{split}
\end{equation*}
Hence
\begin{equation*}
\begin{split}
    &  D^3 F\Big|_{(p',-1)} ((p', 0), (p', 0), (p', 0)) \\
    = & -3 D^2 F\Big|_{(p',-1)} ((p', 0), (p', 0)) + D^3 F\Big|_{(p',-1)} ((0, 1), (0, 1), (0, 1)).
\end{split}
\end{equation*}

To conclude, we rewrite results from Lemma \ref{lem:2.1} to \ref{lem:2.3} with respect to the basis elements
\eqref{ONB} $\{\phi^1, \cdots, \phi^n\}$ as follows:
for $1 \leq \tilde{\alpha}, \tilde{\beta}\leq n$, it holds that
\begin{equation} \label{eq:tau_an}
    \tau^{\tilde{\alpha}n} = \tau^{n \tilde{\alpha} } = |p| \  D^2 F\Big|_{(p',-1)}(\phi^{\tilde{\alpha}}, \phi^n) \sim O \left( \frac{1}{|p|^2}    \right),
\end{equation}
\begin{equation}  \label{eq:D3F_nab}
    D^3 F\Big|_{(p',-1)} (\phi^n,\phi^{\tilde{\alpha}}, \phi^{\tilde{\beta}} ) =  \frac{1}{|p'|} \left( -D^2 F\Big|_{(p',-1)}(\phi^{\tilde{\alpha}}, \phi^{\tilde{\beta}}) + D^3 F\Big|_{(p',-1)}((0,1),\phi^{\tilde{\alpha}}, \phi^{\tilde{\beta}}) \right).
\end{equation}
In particular, we have
\begin{equation}  \label{eq:D3F_nnn}
    D^3 F\Big|_{(p',-1)} (\phi^n,\phi^n, \phi^n ) = \frac{1}{|p'|} \left( -3 D^2 F\Big|_{(p',-1)}(\phi^{n}, \phi^{n}) + \frac{1}{|p'|^2}  D^3 F\Big|_{(p',-1)}((0,1),(0,1),(0,1)) \right).
\end{equation}
From \eqref{eq:D3F_nab} and \eqref{eq:D3F_nnn}, we compute
\begin{eqnarray*}
T^3 (\phi^{\alpha}, \phi^{\alpha}, \phi^{n})    
     & = & - |p| D^3 F\Big|_{(p',-1)} (\phi^{\alpha}, \phi^{\alpha}, \phi^{n}) - \frac{|p'|}{|p|} D^2 F\Big|_{(p',-1)} (\phi^{\alpha}, \phi^{\alpha})    \\
    & = & - \frac{|p|}{|p'|} D^3 F\Big|_{(p',-1)} ((0,1),\phi^{\alpha}, \phi^{\alpha}) + \frac{1}{|p|^2|p'|}  \tau^{\alpha \alpha} \ , \\
    T^3 (\phi^{\alpha}, \phi^n, \phi^n) & = & - |p| D^3 F\Big|_{(p',-1)} (\phi^{\alpha}, \phi^n, \phi^n) - \frac{|p'|}{|p|} D^2 F\Big|_{(p',-1)} (\phi^{\alpha}, \phi^n) \\
     & = & - \frac{|p|}{|p'|} D^3 F\Big|_{(p',-1)} ((0,1), \phi^n, \phi^{\alpha}) + \frac{1}{|p|^2|p'|} \tau^{\alpha n} \ , \\
    T^3 (\phi^{n}, \phi^{n}, \phi^{\alpha}) & = & - |p| D^3 F\Big|_{(p',-1)} (\phi^n, \phi^n, \phi^{\alpha}) \\
    & = & - \frac{|p|}{|p'|} D^3 F\Big|_{(p',-1)} ((0,1), \phi^n, \phi^{\alpha}) + \frac{1}{|p'|} \tau^{\alpha n} \ , \\ 
    T^3 (\phi^{\alpha}, \phi^n, \phi^{\beta}) & = & -|p| D^3 F\Big|_{(p',-1)} (\phi^{\alpha}, \phi^n, \phi^{\beta})  \\
    & = &  - \frac{|p|}{|p'|} D^3 F\Big|_{(p',-1)} ((0,1),\phi^{\alpha}, \phi^{\beta}) + \frac{1}{|p'|} \tau^{\alpha \beta} \ , \\
   T^3 (\phi^{n}, \phi^{n}, \phi^n) & = &  -|p| D^3 F\Big|_{(p',-1)} (\phi^n, \phi^n, \phi^n ) - \frac{|p'|}{|p|} D^2 F\Big|_{(p',-1)} (\phi^n, \phi^n ) \\
   & = & \frac{|p|}{|p'|^3} D^3 F\Big|_{(p',-1)} ((0,1)), (0,1), (0,1)) + \frac{2|p|^2 + 1}{|p'| |p|^2} \tau^{n n} \ . 
\end{eqnarray*}
Then Lemmas \ref{lem:2.2} and \ref{lem:2.3}, estimate \eqref{eq:tau_an} and the
homogeneity of $F$ suggest that there exists a constant $C_4$ such that 
\begin{equation*}
    |T^3 (\phi^{\alpha}, \phi^{\alpha}, \phi^{\beta})| \leq \frac{C_4}{|p|},
\end{equation*}
\begin{equation*}
     |T^3 (\phi^{\alpha}, \phi^{\alpha}, \phi^{n})|,\ |T^3 (\phi^{\alpha}, \phi^n, \phi^n)|, \ |T^3 (\phi^{n}, \phi^{n}, \phi^{\alpha})| \leq \frac{C_4}{|p|^3},
\end{equation*}
\begin{equation*}
    | T^3 (\phi^{\alpha}, \phi^n, \phi^{\beta})| \leq \frac{C_4}{|p|} \quad (\alpha\neq \beta).
\end{equation*}
By the fact that $F$ is convex and homogeneous of degree one, we have that 
$ \tau^{\alpha \alpha}\geq C$ and $\tau^{n n}\geq \frac{C}{|p|^2|}$ for some positive constant $C$.
Hence there are constants $C_5$ and $C_6$ such that
\begin{equation*}
    \frac{C_5}{|p|} \leq T^3 (\phi^{\alpha}, \phi^n, \phi^{\alpha}) \leq \frac{C_6}{|p|},
\end{equation*}
\begin{equation*}
    \frac{C_5}{|p|^3} \leq T^3 (\phi^{n}, \phi^{n}, \phi^n)  \leq \frac{C_6}{|p|^3}
\end{equation*}
completing the proof.
\end{proof}

\end{document}
